\begin{remark}\label{sga3:XXI:2.2.4}
The roots \(\alpha\) and \(\beta\) are orthogonal if and only if the
coroots \(\alpha^\ast\) and \(\beta^\ast\) are orthogonal.
\end{remark}

\begin{lemma}\label{sga3:XXI:2.2.5}
If \(\alpha\) and \(\beta\) are two orthogonal roots, then
\(\alpha+\beta\in R\) if and only if \(\alpha-\beta\in R\).
\end{lemma}

\begin{proof}
Indeed,
\[
 s_\beta(\alpha+\beta)
 =s_\beta(\alpha)+s_\beta(\beta)=\alpha-\beta.
\]
\end{proof}

\begin{lemma}\label{sga3:XXI:2.2.6}
Let \(\alpha\) and \(\beta\) be two nonorthogonal roots. For a coroot
\(\gamma^\ast\), define \(\ell(\gamma^\ast)\) to be the value of
\(\ell\) on the root \(\gamma^\ast\) of \(\cR^\ast\). Then
\begin{equation}
 \ell(\alpha)\ell(\alpha^\ast)
 =\ell(\beta)\ell(\beta^\ast).
 \tag{12}\label{sga3:XXI:eq:12}
\end{equation}
\end{lemma}

\begin{proof}
Multiply formula \eqref{sga3:XXI:eq:11} by the corresponding formula
for \(\cR^\ast\), and use
\((\gamma^\ast)^\ast=\gamma\) for every \(\gamma\in R\). This gives
\[
 (\beta^\ast,\alpha)(\alpha^\ast,\beta)
 \ell(\alpha)\ell(\alpha^\ast)
 =
 (\beta^\ast,\alpha)(\alpha^\ast,\beta)
 \ell(\beta)\ell(\beta^\ast).
\]
The common integer factor is nonzero because the roots are not
orthogonal, so it may be cancelled.
\end{proof}

\subsection{General case}

\begin{proposition}\label{sga3:XXI:2.3.1}
If \(\alpha\) and \(\beta\) are any two roots, then
\begin{equation}
 0\leq(\alpha^\ast,\beta)(\beta^\ast,\alpha)\leq4.
 \tag{13}\label{sga3:XXI:eq:13}
\end{equation}
If \(\alpha\) and \(\beta\) are neither proportional nor orthogonal,
then
\[
 1\leq(\alpha^\ast,\beta)(\beta^\ast,\alpha)\leq3.
\]
\end{proposition}

\begin{proof}
One has
\[
 4(p(\alpha),\beta)^2
 =\ell(\alpha)\ell(\beta)
   (\alpha^\ast,\beta)(\beta^\ast,\alpha).
\]
On the other hand, by Corollary~\ref{sga3:XXI:1.2.6}, the symmetric
bilinear form \((p(x),y)\) is positive and nondegenerate on
\(\VR(R)\). Hence
\[
 (p(x),y)^2\leq\ell(x)\ell(y).
\]
The first inequality follows. In the nonproportional, nonorthogonal
case, the product is a nonzero integer and cannot equal \(4\), which
gives the second inequality.
\end{proof}

\begin{corollary}\label{sga3:XXI:2.3.2}
Let \(\alpha\) and \(\beta\) be two nonorthogonal roots. If
\(\ell(\alpha)=\ell(\beta)\), then there is a \(w\in W\) such that
\(w(\alpha)=\beta\).
\end{corollary}

\begin{proof}
If \(\alpha\) and \(\beta\) are proportional, equality of their
\(\ell\)-values gives \(\alpha=\beta\) or \(\alpha=-\beta\); take
\(w=1\) or \(w=s_\alpha\), respectively. If they are neither
proportional nor orthogonal, formula \eqref{sga3:XXI:eq:11} and
Proposition~\ref{sga3:XXI:2.3.1} give
\[
 (\beta^\ast,\alpha)=(\alpha^\ast,\beta)=\pm1.
\]
If both integers are \(1\), take \(w=s_\beta s_\alpha s_\beta\). If
both are \(-1\), take \(w=s_\alpha s_\beta\).
\end{proof}

\begin{corollary}\label{sga3:XXI:2.3.3}
Let \(\alpha\) and \(\beta\) be roots. If
\(\alpha\neq\beta\) and \((\beta^\ast,\alpha)>0\), then
\(\alpha-\beta\) is a root. Respectively, if
\(\alpha\neq-\beta\) and \((\beta^\ast,\alpha)<0\), then
\(\alpha+\beta\) is a root.
\end{corollary}

\begin{proof}
The second case follows from the first by replacing \(\beta\) by
\(-\beta\). Suppose first that \(\alpha\) and \(\beta\) are
proportional and that \((\beta^\ast,\alpha)>0\). Then
\[
 \beta=\alpha,\qquad 2\beta=\alpha,\qquad\text{or}\qquad
 \beta=2\alpha.
\]
The first case is excluded. In the other two cases,
\(\alpha-\beta=\beta\in R\) and
\(\alpha-\beta=-\alpha\in R\), respectively.

If \(\alpha\) and \(\beta\) are not proportional, then
\((\alpha^\ast,\beta)\) and \((\beta^\ast,\alpha)\) are strictly
positive integers whose product is at most \(3\). Thus one of them
equals \(1\). If \((\beta^\ast,\alpha)=1\), then
\[
 \alpha-\beta=s_\beta(\alpha)\in R;
\]
if \((\alpha^\ast,\beta)=1\), then
\[
 \alpha-\beta=-s_\alpha(\beta)\in R.
\]
\end{proof}

\begin{lemma}\label{sga3:XXI:2.3.4}
Let \(\alpha\) and \(\beta\) be two nonproportional roots. If
\(\beta-\alpha\) is not a root, then
\(\beta+k\alpha\in R\) for \(k=-(\alpha^\ast,\beta)\), but not for
\(k=-(\alpha^\ast,\beta)+1\).
\end{lemma}

\begin{proof}
Indeed,
\[
 \beta-(\alpha^\ast,\beta)\alpha=s_\alpha(\beta)\in R,
\]
whereas
\[
 \beta+\bigl(-(\alpha^\ast,\beta)+1\bigr)\alpha
 =s_\alpha(\beta-\alpha)\notin R.
\]
\end{proof}

\begin{proposition}\label{sga3:XXI:2.3.5}
Let \(\alpha\) and \(\beta\) be two nonproportional roots. The set of
integers \(k\) for which \(\beta+k\alpha\in R\) is an interval
\([p,q]\), with \(p\leq0\) and \(q\geq0\), and
\[
 p+q=-(\alpha^\ast,\beta).
\]
\end{proposition}

\begin{proof}
For the first assertion it suffices, for example, to show that if
\(\beta+k\alpha\in R\) and \(k>0\) is an integer, then
\(\beta+(k-1)\alpha\in R\). This is trivial if \(k=1\). If
\(k\geq2\), then
\[
 (\alpha^\ast,\beta+k\alpha)
 =(\alpha^\ast,\beta)+2k\geq-3+4>0,
\]
and the assertion follows from Corollary~\ref{sga3:XXI:2.3.3}.

Let \([p,q]\) be the interval in question. Applying
Lemma~\ref{sga3:XXI:2.3.4} to \(\beta+p\alpha\), one obtains
\[
 q-p
 =-(\alpha^\ast,\beta+p\alpha)
 =-(\alpha^\ast,\beta)-2p.
\]
The asserted formula follows.
\end{proof}

\begin{remark}\label{sga3:XXI:2.3.6}
The preceding formula contains the qualitative statements of
Lemma~\ref{sga3:XXI:2.2.5} and
Corollary~\ref{sga3:XXI:2.3.3}.
\end{remark}

\begin{complements}\label{sga3:XXI:2.3.7}
\footnote{Editors' note: these complements, useful for the proof of
Proposition~\ref{sga3:XXI:3.1.5}, were added by the editors.}
By formulas \eqref{sga3:XXI:eq:9} and
\eqref{sga3:XXI:eq:10}, and by
Corollary~\ref{sga3:XXI:1.2.6}, the bilinear form on \(\VR(R)\)
defined by
\[
 \langle x,y\rangle=(p(x),y)
\]
is symmetric and positive definite. For every \(\alpha\in R\) and
\(y\in\VR(R)\),
\[
 \langle\alpha,y\rangle
 =\frac{\ell(\alpha)}2(\alpha^\ast,y),
 \qquad
 \ell(\alpha)=\langle\alpha,\alpha\rangle.
\]
Consequently, Corollary~\ref{sga3:XXI:2.3.3} gives the following
corollary.

\medskip
\noindent\textbf{Corollary.}\quad
\SGARefTarget{sga3:XXI:2.3.7:corollary}
Let \(\alpha\neq\beta\) in \(R\). If \(\alpha-\beta\notin R\), then
\[
 \langle\alpha,\beta\rangle\leq0.
\]
\end{complements}

\section{Simple roots, positive roots}

\subsection{Systems of simple roots}

\begin{lemma}\label{sga3:XXI:3.1.1}
Let \(\alpha\) and \(\alpha_i\), \(i=1,\ldots,n\), be roots, and
suppose that \(\alpha\) is distinct from the \(\alpha_i\). If
\[
 \alpha=\sum_{i=1}^n q_i\alpha_i,
 \qquad q_i\in\Qp,
\]
then there is an index \(i\) such that
\[
 q_i\neq0,\qquad(\alpha^\ast,\alpha_i)>0,\qquad
 \alpha-\alpha_i\in R.
\]
\end{lemma}

\begin{proof}
Write
\[
 2=(\alpha^\ast,\alpha)
   =\sum_{i=1}^n q_i(\alpha^\ast,\alpha_i).
\]
This proves the first two assertions. The third then follows from
Corollary~\ref{sga3:XXI:2.3.3}.
\end{proof}
